\documentclass[11pt]{article}
\usepackage[localtheorems]{raeez-math-template}

\newtheorem{theorem}{Theorem}[section]
\newtheorem{proposition}[theorem]{Proposition}
\newtheorem{lemma}[theorem]{Lemma}
\newtheorem{corollary}[theorem]{Corollary}
\theoremstyle{definition}
\newtheorem{definition}[theorem]{Definition}
\theoremstyle{remark}
\newtheorem{remark}[theorem]{Remark}

\DeclareMathOperator{\CoHA}{CoHA}
\DeclareMathOperator{\End}{End}
\DeclareMathOperator{\Hilb}{Hilb}
\DeclareMathOperator{\Rep}{Rep}
\DeclareMathOperator{\Sym}{Sym}
\DeclareMathOperator{\Tr}{Tr}

\newcommand{\Coh}{\mathrm{Coh}}
\newcommand{\Eone}{\mathbb{E}_1}
\newcommand{\Etwo}{\mathbb{E}_2}
\newcommand{\Fock}{\mathcal{F}_{K3}}
\newcommand{\gDel}{\mathfrak{g}_{\Delta_5}}
\newcommand{\Hmuk}{\mathcal{H}_{\mathrm{Muk}}}
\newcommand{\HD}{\mathbf{H}_{\Delta_5}^{\mathrm{cand}}}
\newcommand{\Lmuk}{\widetilde{\Lambda}_{K3}}
\newcommand{\LmukC}{\Lmuk \otimes_{\mathbb{Z}} \mathbb{C}}
\newcommand{\YHeis}{Y^{\mathrm{Heis}}_{\mathrm{Muk}}}

\title{The Mukai Heisenberg Branch of the K3 Yangian}
\author{Raeez Lorgat}
\date{April 2026}

\begin{document}
\maketitle

\begin{abstract}
The abelian K3 Yangian presentation is the Mukai-Heisenberg branch:
a rank-$24$ abelian current algebra with diagonal RTT/free-field
quantisation. Its input is the Mukai lattice
$\Lmuk \simeq \mathrm{II}_{4,20}$, a Miura frame, and parameters
$h_1,\ldots,h_{24}$ with vanishing first Newton sum. The construction
is abelian at Lie level. Vertex-algebra, quantum-toroidal, and
Borcherds structures require additional input. The Hall-Drinfeld branch
for $K3 \times E$ is a separate $d=3$ Borcherds construction. The
integer $324$ is the coefficient $\chi(\Hilb^2(K3))$ in the
Nakajima-Grojnowski Fock character; finite module counts require a
separate tensor category.
\end{abstract}

\tableofcontents

\section{Mukai Input}
\label{sec:mukai-input}

Let $S$ be a complex projective K3 surface. Its Mukai lattice is
\[
  \Lmuk
  =
  H^0(S,\mathbb{Z}) \oplus H^2(S,\mathbb{Z}) \oplus H^4(S,\mathbb{Z}),
\]
with pairing
\[
  \bigl((r,D,s),(r',D',s')\bigr)_{\mathrm{Muk}}
  =
  D \cdot D' - rs' - r's .
\]
By Mukai, $\Lmuk$ is even, unimodular, and has signature $(4,20)$.
The second-cohomology lattice $H^2(S,\mathbb{Z})$ has rank $22$ and
signature $(3,19)$. The rank controlling the abelian K3 Yangian is the
Mukai rank $24$.

\begin{definition}[Miura frame]
\label{def:miura-frame}
A Miura frame is a complex orthogonal basis
$e_1,\ldots,e_{24}$ of $\LmukC$ for which
\[
  (e_i,e_j)_{\mathrm{Muk}} = \varepsilon_i \delta_{ij},
  \qquad
  \varepsilon_i =
  \begin{cases}
    +1, & 1 \leq i \leq 4,\\
    -1, & 5 \leq i \leq 24,
  \end{cases}
\]
together with parameters $h_1,\ldots,h_{24} \in \mathbb{C}$ satisfying
\[
  p_1(h) := \sum_{i=1}^{24} h_i = 0 .
\]
The equation $p_1(h)=0$ is the vanishing of the first Newton sum in
the chosen Miura frame. It is the parameter-level trace condition; the
signature $(4,20)$ is carried by the bilinear form, not by this scalar
equation.
\end{definition}

\begin{definition}[Mukai-Heisenberg current algebra]
\label{def:mukai-heisenberg}
The Mukai-Heisenberg current algebra $\Hmuk$ is generated by modes
$J_{i,n}$, $1 \leq i \leq 24$, $n \in \mathbb{Z}$, and a central unit
$\mathbf{1}$, with relations
\[
  [J_{i,m},J_{j,n}]
  =
  m\,\varepsilon_i\,\delta_{ij}\,\delta_{m+n,0}\,\mathbf{1}.
\]
The current
\[
  J_i(z) = \sum_{n \in \mathbb{Z}} J_{i,n} z^{-n-1}
\]
has operator-product expansion
\[
  J_i(z)J_j(w) \sim
  \frac{\varepsilon_i\delta_{ij}}{(z-w)^2}.
\]
The Fock representation is
\[
  \Fock =
  \Sym\!\left(
    \bigoplus_{n\geq 1}\LmukC\, t^{-n}
  \right),
\]
where $J_{i,-n}$ acts by multiplication and $J_{i,n}$ acts by
$n\varepsilon_i\,\partial/\partial J_{i,-n}$ for $n>0$.
\end{definition}

\begin{definition}[The abelian K3 Yangian branch]
\label{def:heisenberg-branch}
The algebra denoted here by $\YHeis$ is the diagonal RTT/free-field
quantisation of $\Hmuk$ determined by the Miura polynomial
\[
  T_{K3}(u;z)
  =
  :\prod_{i=1}^{24}\bigl(u-\phi_i(z)\bigr):
  =
  \sum_{s=0}^{24}(-1)^s\psi_s(z)u^{24-s},
\]
where the $\phi_i$ are the free fields of Definition
\ref{def:mukai-heisenberg}. Its scalar structure function is
\[
  g_{K3}(u)
  =
  \prod_{i=1}^{24}\frac{u-h_i}{u+h_i}.
\]
This is the abelian Lie-current branch. The word ``Yangian'' in this
notation refers to the diagonal RTT/free-field package, not to a
Drinfeld $J$-presentation for a non-abelian finite-dimensional simple
Lie algebra.
\end{definition}

\section{Algebraic Content}
\label{sec:algebraic-content}

\begin{proposition}[PBW, Fock module, and tensor factorisation]
\label{prop:pbw-fock-factorisation}
The Mukai-Heisenberg branch has the following properties.
\begin{enumerate}[label=\textup{(\alph*)}]
\item The enveloping algebra of $\Hmuk$ has Poincare-Birkhoff-Witt
basis ordered by negative modes, zero modes, and positive modes.
\item The Fock representation $\Fock$ is generated from a vacuum
$|0\rangle$ annihilated by all $J_{i,n}$ with $n>0$.
\item In a Miura frame,
\[
  \Hmuk \cong \widehat{\bigoplus}_{i=1}^{24}\mathrm{Heis}_{\varepsilon_i},
  \qquad
  \YHeis \cong
  \widehat{\bigotimes}_{i=1}^{24}Y^{\mathrm{Heis}}_i
\]
as completed current algebras with diagonal RTT factors.
\end{enumerate}
The factor $Y^{\mathrm{Heis}}_i$ is the rank-one Heisenberg current
factor. Identifying it with a specialisation of
$Y(\widehat{\mathfrak{gl}}_1)$ requires an additional affine-Yangian
normalisation and kills the generic cubic shuffle interaction.
\end{proposition}

\begin{proof}
The PBW statement is the standard PBW theorem for a central extension
of an abelian loop algebra. The Fock action follows from
\[
  \left[
    n\varepsilon_i\frac{\partial}{\partial J_{i,-n}},
    J_{j,-m}
  \right]
  =
  n\varepsilon_i\delta_{ij}\delta_{nm}.
\]
Since the Mukai form is diagonal in the chosen frame, different
directions commute. The completed tensor product is therefore the
universal algebra generated by the $24$ rank-one factors with pairwise
commuting images.
\end{proof}

\begin{proposition}[Structure-function identities]
\label{prop:structure-function}
The scalar structure function satisfies
\[
  g_{K3}(u)g_{K3}(-u)=1
\]
and
\[
  \log g_{K3}(u)
  =
  -2\sum_{r\geq 0}
  \frac{p_{2r+1}(h)}{(2r+1)u^{2r+1}},
  \qquad
  p_k(h)=\sum_{i=1}^{24}h_i^k .
\]
Thus the $u^{-1}$ coefficient vanishes exactly when $p_1(h)=0$.
The Mukai signature enters the two-point function through the signs
$\varepsilon_i$, and the total Heisenberg level is
\[
  \Tr(\varepsilon)=4-20=-16 .
\]
\end{proposition}

\begin{proof}
The identity $g(u)g(-u)=1$ follows by pairing each factor
$(u-h_i)/(u+h_i)$ with $(-u-h_i)/(-u+h_i)$. The logarithmic expansion is
the difference of the two Taylor series for
$\log(1-h_i/u)$ and $\log(1+h_i/u)$, summed over $i$.
The trace is the difference between the positive and negative
dimensions of the Mukai form.
\end{proof}

\begin{remark}[Kummer frame]
\label{rem:kummer-frame}
The common Kummer frame
\[
  h_i=\varepsilon \quad (1\leq i\leq 4),
  \qquad
  h_i=-\varepsilon/5 \quad (5\leq i\leq 24)
\]
satisfies $p_1(h)=4\varepsilon-20\varepsilon/5=0$. In this frame
$g_{K3}$ has four zeros at $u=\varepsilon$, twenty zeros at
$u=-\varepsilon/5$, four poles at $u=-\varepsilon$, and twenty poles
at $u=\varepsilon/5$. These multiplicities display the signature block
sizes. Serre relations require enhanced root-algebra data.
\end{remark}

\section{Serre Relations}
\label{sec:serre-relations}

\begin{proposition}[Serre content of the abelian branch]
\label{prop:abelian-serre}
At the Lie-current level the Mukai-Heisenberg branch has no
non-trivial Serre relations. For $i\neq j$,
\[
  [J_i(z),J_j(w)] = 0,
\]
and every iterated adjoint expression built from distinct Heisenberg
directions vanishes for the elementary reason that the Lie algebra is
abelian modulo its central extension.
\end{proposition}

\begin{proof}
The defining commutator in Definition \ref{def:mukai-heisenberg} is
central and is zero unless $i=j$ and the mode numbers sum to zero.
Consequently
\[
  \operatorname{ad}(J_{i,m})\operatorname{ad}(J_{j,n})(x)=0
\]
for all $i\neq j$ and all $x\in\Hmuk$. This is plain Heisenberg
commutativity. Kac-Moody Serre cancellation belongs to enhanced root
data.
\end{proof}

\begin{remark}[Enhancement and Borcherds sectors]
\label{rem:enhancement-borcherds}
ADE Serre relations at an $E_8\oplus E_8$ enhancement point are
relations of the enhanced root algebra on that special locus. The
Borcherds real-root and imaginary-root relations of
$\gDel$ are relations of the BKM denominator algebra attached to
$\Delta_5$. Neither family is a relation of $\YHeis$ at a generic
K3 period point.
\end{remark}

\section{Fock Character and Hilbert Schemes}
\label{sec:fock-character}

\begin{theorem}[Nakajima-Grojnowski realisation]
\label{thm:nakajima-grojnowski}
There is an action of the Heisenberg algebra
$H^*(S,\mathbb{C})\otimes t\mathbb{C}[t,t^{-1}]$ on
\[
  \mathcal{V}_S :=
  \bigoplus_{n\geq 0} H^*(\Hilb^n S,\mathbb{C})
\]
for which $\mathcal{V}_S$ is the Fock representation generated by
$H^*(S,\mathbb{C})$. The Euler-characteristic generating series is
\[
  \sum_{n\geq 0}\chi(\Hilb^n S)q^n
  =
  \prod_{m\geq 1}(1-q^m)^{-24}
  =
  1+24q+324q^2+3200q^3+\cdots .
\]
Equivalently,
\[
  \eta(\tau)^{-24}
  =
  q^{-1}+24+324q+3200q^2+\cdots .
\]
\end{theorem}

\begin{proof}
This is the Nakajima-Grojnowski Heisenberg construction for Hilbert
schemes of points on a smooth projective surface, specialised to K3.
G\"ottsche's formula gives the Euler-characteristic series. The
coefficient of $q^2$ in the product is
\[
  \binom{24+1}{2}+24=300+24=324,
\]
coming respectively from two degree-one creation operators and one
degree-two creation operator.
\end{proof}

\begin{corollary}[Meaning of the integer $324$]
\label{cor:meaning-324}
The integer $324$ is
\[
  324=\chi(\Hilb^2 S)=p_{24}(2),
\]
the second coefficient of the $24$-boson Fock character. Module counts
for the abelian Heisenberg algebra $\Hmuk$ or for $\YHeis$ require
additional finite tensor-categorical data.
\end{corollary}

\begin{proof}
The Heisenberg algebra has Fock modules parametrised by highest
weights. A finite list of $324$ irreducibles appears only after
supplying a separate finite tensor category, such as a root-of-unity
small quantum group or a finite quotient. The Nakajima-Grojnowski
theorem supplies a graded vector-space coefficient. Such a tensor
category is additional data.
\end{proof}

\section{CoHA Boundary}
\label{sec:coha-boundary}

\begin{proposition}[What the Hilbert-scheme theorem proves]
\label{prop:coha-boundary}
The theorem of Nakajima and Grojnowski proves a Heisenberg action on
$\bigoplus_n H^*(\Hilb^n S)$. An algebra isomorphism
\[
  \CoHA(S) \cong \YHeis
\]
or an isomorphism between any compact K3 Hall algebra and the diagonal
RTT algebra of Definition \ref{def:heisenberg-branch} requires a
specified Hall product and an intertwining theorem.
\end{proposition}

\begin{proof}
The Heisenberg construction uses correspondences between Hilbert
schemes and produces creation and annihilation operators. A Hall
algebra comparison requires a specified moduli stack, a product, a
coproduct or pairing when a double is intended, and a proof that the
Nakajima operators intertwine those structures. Equality of Fock
characters is insufficient: different algebra structures can act on
the same graded vector space.
\end{proof}

\begin{remark}[Positive halves and doubles]
\label{rem:positive-halves-doubles}
For local surfaces and quiver examples, cohomological Hall algebras
often realise positive halves of quantum groups. Passing from a
positive half to a full quantum group requires a Hopf pairing and a
Drinfeld double. This distinction is load-bearing for the K3 and
$K3\times E$ constructions.
\end{remark}

\section{Separation from the Hall-Drinfeld BKM Branch}
\label{sec:hall-drinfeld-separation}

The K3 Yangian branch and the $K3\times E$ BKM branch share the Mukai
lattice but arise from different inputs.

\begin{proposition}[Two inputs]
\label{prop:two-inputs}
The abelian K3 Yangian branch is the $d=2$ Mukai-Heisenberg
construction attached to $D^b\Coh(S)$. The Hall-Drinfeld BKM branch is
a $d=3$ construction problem attached to $D^b\Coh(S\times E)$ and to
the Borcherds denominator $\Delta_5$.
\end{proposition}

\begin{proof}
For $d=2$, the native chiral output is $\Etwo$-chiral and the
Mukai-Heisenberg currents give the abelian Lie-current stratum. For
$d=3$, the native output is $\Eone$-chiral. A braided or $\Etwo$
structure appears only after passing to a centre, double, or an
additional stage-$2$ specialisation under hypotheses. Thus the
Hall-Drinfeld object belongs to the centre/double side of the
$d=3$ theory, not to the native abelian K3 Yangian.
\end{proof}

\begin{definition}[BKM Hall-Drinfeld target]
\label{def:bkm-hall-drinfeld-target}
The $K3\times E$ comparison target has the form
\[
  \HD
  :=
  \mathcal{D}_{\hbar}\!\left(
    Y^+_{\mathrm{Hall}}(K3\times E),\,
    D_5(2Z)=64^{-1}\Delta_5(2Z),\,
    R_{\mathrm{Sieg,dyn}}
  \right),
\]
where $Y^+_{\mathrm{Hall}}(K3\times E)$ is a positive Hall half,
$D_5$ is the Borcherds denominator input, and
$R_{\mathrm{Sieg,dyn}}$ is the Siegel-elliptic dynamical
$R$-matrix datum. This is a Hall-Drinfeld BKM target. Yangian
structures belong to the separate Mukai-Heisenberg branch.
\end{definition}

\begin{remark}[The four invariants of $K3\times E$]
\label{rem:four-invariants}
The canonical spectrum attached to $K3\times E$ is
\[
  \{
    \kappa_{\mathrm{cat}},
    \kappa_{\mathrm{ch}}^{\mathrm{Heis}},
    \kappa_{\mathrm{BKM}},
    \kappa_{\mathrm{fiber}}
  \}(K3\times E)
  =
  \{0,3,5,24\}.
\]
Here
\[
  \kappa_{\mathrm{cat}}(K3\times E)
  =
  \chi(\mathcal{O}_{K3})\chi(\mathcal{O}_E)
  =
  2\cdot 0=0,
\]
$\kappa_{\mathrm{ch}}^{\mathrm{Heis}}=3$ belongs to the
Heisenberg-Mukai stage-$2$ specialisation,
$\kappa_{\mathrm{BKM}}(\Delta_5)=c_1(0)/2=5$ is the Borcherds weight,
and $\kappa_{\mathrm{fiber}}=24$ is the Mukai-lattice rank of the K3
fibre. The value $2$ is $\kappa_{\mathrm{cat}}(K3)$ for the fibre
alone. The total-space invariant is $0$.
\end{remark}

\section{Root of Unity and Modular Data}
\label{sec:root-of-unity}

\begin{proposition}[Finite modular data require extra input]
\label{prop:no-abelian-324-module-theorem}
A finite modular tensor category with $324$ simple objects requires
additional input beyond the abelian branch $\YHeis$. In a
positive-definite lattice-VOA formula the number of simple currents is
the order of the discriminant group $\Lambda^\vee/\Lambda$; for the
Mukai lattice this group is trivial because $\Lmuk$ is unimodular. For
the indefinite Mukai lattice, rational modular tensor categories require
extra analytic or categorical input.
\end{proposition}

\begin{proof}
Unimodularity gives $\Lmuk^\vee=\Lmuk$. Hence the standard finite
discriminant-form modular datum has one simple-current label. The
Heisenberg VOA itself has continuous Fock labels before a lattice or
finite quotient is imposed. These constructions produce no $324$ simple
objects.
\end{proof}

\begin{remark}[Where root-of-unity counts live]
\label{rem:root-of-unity-counts}
A statement that a root-of-unity small quantum group attached to
$\gDel$ has $324$ simple modules at $N=2$ belongs to the BKM
Hall-Drinfeld branch. Its proof requires an integral form, a
Kazhdan-Lusztig negligible quotient, a specified tensor category, and
an identification of its simples with the coefficient
$p_{24}(2)=324$. The Fock-character coefficient supplies the number;
the finite category is extra structure.
\end{remark}

\section{Abelian Theorem}
\label{sec:abelian-theorem}

\begin{theorem}[Scope of the abelian K3 Yangian presentation]
\label{thm:abelian-k3-yangian-scope}
Let $S$ be a K3 surface, let $\Lmuk$ be its Mukai lattice, choose a
Miura frame, and let $h_1,\ldots,h_{24}$ satisfy $\sum_i h_i=0$.
Then the construction above gives a well-defined rank-$24$
Mukai-Heisenberg branch $\YHeis$ with:
\begin{enumerate}[label=\textup{(\alph*)}]
\item modes $J_{i,n}$ satisfying the central Heisenberg relation
\[
  [J_{i,m},J_{j,n}]
  =
  m\varepsilon_i\delta_{ij}\delta_{m+n,0}\mathbf{1};
\]
\item Fock module
$\Fock=\Sym(\bigoplus_{n\geq 1}\LmukC t^{-n})$;
\item Miura transfer polynomial
$T_{K3}(u;z)=:\prod_i(u-\phi_i(z)):$ of degree $24$;
\item scalar structure function
$g_{K3}(u)=\prod_i(u-h_i)/(u+h_i)$ satisfying
$g_{K3}(u)g_{K3}(-u)=1$;
\item vanishing $u^{-1}$ coefficient in $\log g_{K3}(u)$ exactly when
$\sum_i h_i=0$;
\item tensor factorisation into $24$ commuting rank-one Heisenberg
current factors;
\item Hilbert-scheme Fock character
\[
  \prod_{m\geq 1}(1-q^m)^{-24}
  =
  1+24q+324q^2+\cdots .
\]
\end{enumerate}
These are abelian Lie-current and free-field statements.
\end{theorem}

\begin{proof}
Items (a), (b), and (f) are Definition \ref{def:mukai-heisenberg} and
Proposition \ref{prop:pbw-fock-factorisation}. Items (c), (d), and
(e) are Definition \ref{def:heisenberg-branch} and Proposition
\ref{prop:structure-function}. Item (g) is Theorem
\ref{thm:nakajima-grojnowski}.
\end{proof}

\begin{corollary}[Proof obligations beyond the abelian theorem]
\label{cor:outside-abelian-theorem}
The abelian theorem proves exactly the statements above. The following
assertions require additional data:
\begin{enumerate}[label=\textup{(\arabic*)}]
\item $\YHeis$ is the BKM algebra $\gDel$ or the Hall-Drinfeld target
$\HD$.
\item $\YHeis$ has non-trivial ADE or Borcherds Serre relations at a
generic K3 period point.
\item $\YHeis$ is isomorphic to a compact K3 cohomological Hall
algebra.
\item The coefficient $324$ is a count of irreducible modules of
$\YHeis$.
\item The $K3\times E$ invariant $\kappa_{\mathrm{cat}}$ equals $2$.
\end{enumerate}
\end{corollary}

\begin{proof}
Item (1) is Proposition \ref{prop:two-inputs} and Definition
\ref{def:bkm-hall-drinfeld-target}. Item (2) is Proposition
\ref{prop:abelian-serre}. Item (3) is Proposition
\ref{prop:coha-boundary}. Item (4) is Corollary
\ref{cor:meaning-324} and Proposition
\ref{prop:no-abelian-324-module-theorem}. Item (5) is Remark
\ref{rem:four-invariants}.
\end{proof}

\section{Proof Obligations Beyond the Abelian Branch}
\label{sec:proof-obligations}

A non-abelian theorem extending the abelian branch must supply the
following data.

\begin{enumerate}[label=\textup{(\arabic*)}]
\item A non-abelian current algebra or BKM input with explicit real
and imaginary root generators, multiplicities, and Borcherds
relations. The Mukai-Heisenberg algebra alone supplies only the
Cartan/free-field sector.
\item A quantisation functor or RTT presentation for that non-abelian
input, including coproduct, antipode, Hopf pairing, and PBW theorem.
\item For the $K3\times E$ branch, a compact Hall positive half
$Y^+_{\mathrm{Hall}}(K3\times E)$, a nondegenerate Hopf pairing, and
the completed Drinfeld double regulated by the $\Delta_5$ denominator.
\item For any root-of-unity module count, an integral form, the
negligible quotient, and a proof that the resulting simple objects are
counted by $p_{24}(2)=324$.
\item For any comparison with Maulik-Okounkov stable envelopes, a
global K3 construction replacing the equivariant toric input available
for $\mathbb{A}^2$ and $\mathbb{C}^3$.
\end{enumerate}

Until these data are supplied, the abelian theorem remains exactly the
Mukai-Heisenberg statement of Theorem
\ref{thm:abelian-k3-yangian-scope}.

\bibliographystyle{plain}
\begin{thebibliography}{99}
\bibitem{Borcherds1988}
R.~Borcherds,
\emph{Generalized Kac-Moody algebras},
J.~Algebra~115 (1988), 501--512.

\bibitem{Drinfeld1985}
V.~Drinfeld,
\emph{Hopf algebras and the quantum Yang-Baxter equation},
Sov.~Math.~Dokl.~32 (1985), 254--258.

\bibitem{Drinfeld1988}
V.~Drinfeld,
\emph{A new realization of Yangians and quantized affine algebras},
Sov.~Math.~Dokl.~36 (1988), 212--216.

\bibitem{Goettsche1990}
L.~G\"ottsche,
\emph{The Betti numbers of the Hilbert scheme of points on a smooth
projective surface},
Math.~Ann.~286 (1990), 193--207.

\bibitem{GritsenkoNikulin1998}
V.~Gritsenko and V.~Nikulin,
\emph{Automorphic forms and Lorentzian Kac-Moody algebras II},
Int.~J.~Math.~9 (1998), 201--275.

\bibitem{Grojnowski1996}
I.~Grojnowski,
\emph{Instantons and affine algebras I: the Hilbert scheme and vertex
operators},
Math.~Res.~Lett.~3 (1996), 275--291.

\bibitem{KazhdanLusztig1993}
D.~Kazhdan and G.~Lusztig,
\emph{Tensor structures arising from affine Lie algebras I},
J.~Amer.~Math.~Soc.~6 (1993), 905--947.

\bibitem{Lusztig1990}
G.~Lusztig,
\emph{Quantum groups at roots of unity},
Geom.~Dedicata~35 (1990), 89--113.

\bibitem{MaulikOkounkov2012}
D.~Maulik and A.~Okounkov,
\emph{Quantum groups and quantum cohomology},
arXiv:1211.1287.

\bibitem{Mukai1987}
S.~Mukai,
\emph{On the moduli space of bundles on K3 surfaces, I},
Tata Inst.\ Fund.\ Res.\ Stud.~Math.~11 (1987).

\bibitem{Nakajima1997}
H.~Nakajima,
\emph{Heisenberg algebra and Hilbert schemes of points on projective
surfaces},
Ann.~Math.~145 (1997), 379--388.

\bibitem{SchiffmannVasserot2013}
O.~Schiffmann and E.~Vasserot,
\emph{Cherednik algebras, W-algebras and the equivariant cohomology
of the moduli space of instantons on $\mathbb{A}^2$},
Publ.~IHES~118 (2013), 213--342.

\bibitem{Tsymbaliuk2017}
A.~Tsymbaliuk,
\emph{The affine Yangian of $\mathfrak{gl}_1$ revisited},
Adv.~Math.~304 (2017), 583--645.
\end{thebibliography}

\end{document}
